%%%% Настройки автора
%%
%%   Пожалуйста, ознакомьтесь с функционалом шаблона из [1,2], а также с пакетами, подключенными в ch_preamble.
%%

\addbibresource{my_folder/my_biblio.bib}
\ExecuteBibliographyOptions{defernumbers=false,sorting=none}

\hypersetup{
pdfstartview={FitBH}
}

\let\ordinal\relax
\usepackage{datetime}
\usepackage{tikz}
\usetikzlibrary{arrows.meta,positioning,shapes.geometric}

\lstdefinelanguage{JavaScript}{
  keywords={const, let, var, function, return, if, else, for, while, break, continue, switch, case, default, try, catch, finally, throw, new, class, extends, import, from, export, async, await, true, false, null, undefined},
  sensitive=true,
  morecomment=[l]{//},
  morecomment=[s]{/*}{*/},
  morestring=[b]',
  morestring=[b]"
}

\newcommand{\overbar}[1]{\mkern 1.5mu\overline{\mkern-1.5mu#1\mkern-1.5mu}\mkern 1.5mu}
\DeclareMathAlphabet{\mathbbmsl}{U}{bbm}{m}{sl}
\newcommand{\cont}[1][K]{\ensuremath{\mathbbmsl{#1}}}
\newcommand{\fcaarrow}[1]{%
	{}^{\scriptscriptstyle\bm{#1}}
}
\newcommand{\uA}{\fcaarrow{{\uparrow\mkern-12mu}}}
\newcommand{\dA}{\fcaarrow{\downarrow\mkern-2mu}}
\newcommand{\ud}{\fcaarrow{\uparrow\mkern-12mu}\fcaarrow{\downarrow\mkern-2mu}}
\newcommand{\du}{\fcaarrow{\downarrow\mkern-2mu}\fcaarrow{\uparrow\mkern-12mu}}
\newcommand\eqdef{\mathrel{\overset{\makebox[0pt]{\mbox{\normalfont\tiny def}}}{=}}}

\theoremstyle{myplain}
\newtheorem{m-new-env-first}{Название\_окружения}[chapter]

\theoremstyle{mydefinition}
\newtheorem{m-new-env-second}{Название\_окружения}[chapter]
